\documentclass[conference]{IEEEtran}
\IEEEoverridecommandlockouts

\usepackage{amsmath,amssymb,amsfonts}
\usepackage{graphicx}
\usepackage{url}
\usepackage{listings}
\usepackage{float}

\lstset{
    language=C++,
    basicstyle=\ttfamily\scriptsize,
    breaklines=true,
    columns=fullflexible,
    frame=single,
    tabsize=2
}

\begin{document}
\raggedbottom

\title{Discrete Fourier Transform and Spatial Convolution Filtering Using OpenCV\\
}

\author{\IEEEauthorblockN{Aufa Rienaldifaza Ahmad - 13524027}
\IEEEauthorblockA{\textit{Sekolah Teknik Elektro dan Informatika} \\
\textit{Institut Teknologi Bandung}\\
Bandung, Indonesia \\
aufa.rienaldifaza.a@gmail.com}
}

\maketitle

\section{Theoretical Foundation}

\subsection{Linear Spatial Filtering}
\label{sec:theory-conv}

Convolution combines every part of an image with a kernel: a fixed
size array of coefficients with an anchor point at its
center~\cite{opencv-filter2d}. The kernel is rotated $180^\circ$
about its anchor, then slid over the image; at each position its
coefficients are multiplied with the underlying pixels, summed, and
written back at the anchor. For an image $I$ and kernel $K$ of size
$M_i \times M_j$ with anchor $(a_i, a_j)$,

\begin{equation}
    H(x, y) = \sum_{i=0}^{M_i - 1} \sum_{j=0}^{M_j - 1}
    K(i, j)\, I(x - i + a_i,\; y - j + a_j)
    \label{eq:convolution}
\end{equation}

\noindent OpenCV's \texttt{cv::filter2D} instead computes
\emph{correlation} (the same sum without the $180^\circ$ rotation), so
a kernel must be pre-rotated to obtain true convolution through it;
kernels symmetric under $180^\circ$ rotation give identical results
either way.

\subsection{Discrete Fourier Transform}
\label{sec:theory-dft}

The Fourier Transform decomposes an image into sine and cosine
components in the frequency domain~\cite{opencv-dft}. For an
$N \times N$ image $f(i,j)$,

\begin{equation}
    F(k, l) = \sum_{i=0}^{N-1} \sum_{j=0}^{N-1}
    f(i, j)\, e^{-i 2\pi \left(\frac{ki}{N} + \frac{lj}{N}\right)}
    \label{eq:dft}
\end{equation}

\noindent which is complex-valued even for real-valued input. Since
\texttt{cv::Mat} has no native complex type, the real and (zero)
imaginary planes are interleaved with \texttt{cv::merge} before the
transform and separated with \texttt{cv::split} after. DFT
performance is best for sizes that are products of small primes, so
the input is zero-padded to \texttt{cv::getOptimalDFTSize} first.

A complex value cannot be shown as a pixel, so the spectrum is
visualized through its magnitude,

\begin{equation}
    M(k, l) = \sqrt{\operatorname{Re}(F(k,l))^2 + \operatorname{Im}(F(k,l))^2}
    \label{eq:magnitude}
\end{equation}

\noindent The DC term dominates $M$ by orders of magnitude, so a
logarithmic transform $M_1 = \log(1 + M)$ is applied before display.
The raw DFT also places the zero-frequency term at the corner
$(0,0)$; swapping diagonally opposite quadrants recenters it for
interpretability.

\section{Experimental Results}
\label{sec:experiments}

\subsection{Implementation}
The methods described in Section~\ref{sec:theory-conv} and
Section~\ref{sec:theory-dft} were implemented as member functions of
an \texttt{ImageProcessor} class, built as an OpenCV-based
\texttt{ament\_cmake} library. Three filtering-related methods were
implemented:

\begin{itemize}
    \item \texttt{convolutionFilter}: applies a user-supplied kernel
    using OpenCV's \texttt{cv::filter2D}, which performs correlation
    rather than true convolution (Section~\ref{sec:theory-conv}), so
    the kernel is first rotated $180^\circ$ with \texttt{cv::rotate}
    so that the combined operation matches Eq.~\ref{eq:convolution}.

    \item \texttt{convolutionFilterManual}: a from-scratch
    implementation of the same mathematical operation. The kernel is applied as a convolution directly,
    by indexing it in reverse order rather than pre-rotating it, and
    out-of-bounds neighbors are treated as zero, matching
    \texttt{convolutionFilter}'s border behavior.

    \item \texttt{fourierTransform}: computes the log-magnitude
    spectrum of an image following the pipeline in
    Section~\ref{sec:theory-dft} (optimal-size padding, complex
    plane construction, DFT, magnitude, log scaling, and optional
    quadrant centering).
\end{itemize}

All three methods were exercised through a demo executable on the
512$\times$512, 3-channel \texttt{Lenna.png} test image, shown
unmodified in Fig.~\ref{fig:lenna_original}.

\begin{figure}[H]
    \centering
    \includegraphics[width=0.85\columnwidth]{assets/lenna_original.png}
    \caption{Original 512$\times$512 test image (\texttt{Lenna.png}).}
    \label{fig:lenna_original}
\end{figure}

\subsection{Correctness of the Manual Convolution}
\texttt{convolutionFilterManual} was validated against
\texttt{convolutionFilter} using a $15\times15$ Gaussian kernel from
\texttt{cv::getGaussianKernel}. Results are shown in Fig.~\ref{fig:lenna_filtered_opencv} and
Fig.~\ref{fig:lenna_filtered_manual}.

\begin{figure}[H]
    \centering
    \includegraphics[width=0.85\columnwidth]{assets/lenna_filtered_opencv.png}
    \caption{Result of \texttt{convolutionFilter} (\texttt{filter2D}-based) using the $15\times15$ Gaussian kernel $K$.}
    \label{fig:lenna_filtered_opencv}
\end{figure}

\begin{figure}[H]
    \centering
    \includegraphics[width=0.85\columnwidth]{assets/lenna_filtered_manual.png}
    \caption{Result of \texttt{convolutionFilterManual} using the same kernel $K$, visually indistinguishable from Fig.~\ref{fig:lenna_filtered_opencv}.}
    \label{fig:lenna_filtered_manual}
\end{figure}

\subsection{Performance Comparison}
\begin{table}[H]
\caption{Average Runtime per Frame (512$\times$512$\times$3 image, $3\times3$ kernel)}
\begin{center}
\begin{tabular}{|l|c|}
\hline
\textbf{Method} & \textbf{Time (ms/frame)} \\
\hline
\texttt{convolutionFilter} (\texttt{filter2D}) & 0.27 \\
\hline
\texttt{convolutionFilterManual} (loops) & 11.96 \\
\hline
\end{tabular}
\label{tab:performance}
\end{center}
\end{table}

\noindent The manual implementation is approximately $45\times$ slower
than OpenCV's \texttt{filter2D}. This is due to \texttt{filter2D}
being backed by vectorized (SIMD), cache-optimized, and potentially
multi-threaded code, whereas the manual version performs one scalar
multiply-accumulate per kernel tap through plain nested loops.

\subsection{Fourier Spectrum}
The log-magnitude spectrum computed by \texttt{fourierTransform}
shows a bright peak at the center
corresponding to the zero-frequency (DC) component, surrounded by a
star-shaped pattern of radial spikes aligned with the image's
row and column edges after border padding. Disabling the quadrant
centering step reproduces the same spectrum with the DC term located
at the corners instead, consistent with the raw output layout of
\texttt{cv::dft}.

\begin{figure}[H]
    \centering
    \includegraphics[width=0.85\columnwidth]{assets/lenna_fft_spectrum.png}
    \caption{Log-magnitude Fourier spectrum of \texttt{Lenna.png}, quadrant-centered so the zero-frequency term appears at the middle of the image.}
    \label{fig:lenna_fft_spectrum}
\end{figure}

\subsection{Effect of Kernel Size on Blur Strength}
\label{sec:kernel-sweep}
Initial testing with a small $3\times3$ box-averaging kernel
$K = \frac{1}{9}\mathbf{1}_{3\times3}$ produced only a mild blur. Box
kernels $K_N = \frac{1}{N^2}\mathbf{1}_{N\times N}$ were therefore
applied to \texttt{Lenna.png} via \texttt{convolutionFilter} for
$N \in \{3, 9, 15, 21, 31\}$, shown in
Fig.~\ref{fig:kernel_size_comparison}. Blur strength increases
visibly and steadily with $N$, becoming clearly soft from around
$N=15$ onward.

\begin{figure}[H]
    \centering
    \includegraphics[width=\columnwidth]{assets/kernel_size_comparison.png}
    \caption{Box-blur result on \texttt{Lenna.png} for increasing kernel size $N$, using \texttt{convolutionFilter}.}
    \label{fig:kernel_size_comparison}
\end{figure}

\onecolumn
\appendix
\section{Source Code}
The full project source is available at
\url{https://github.com/aufafaza/EL4007-Digital-Image-Processing_project}.
The five \texttt{ImageProcessor} methods are reproduced below for
reference.

\noindent\textbf{\texttt{toGrayscale}}
\begin{lstlisting}
cv::Mat ImageProcessor::toGrayscale(const cv::Mat & input) const
{
  if (input.empty()) {
    throw std::invalid_argument("ImageProcessor::toGrayscale: input image is empty");
  }
  if (input.channels() == 1) {
    return input.clone();
  }
  cv::Mat gray;
  cv::cvtColor(input, gray, cv::COLOR_BGR2GRAY);
  return gray;
}
\end{lstlisting}

\noindent\textbf{\texttt{gaussianBlur}}
\begin{lstlisting}
cv::Mat ImageProcessor::gaussianBlur(const cv::Mat & input, int kernel_size) const
{
  if (input.empty()) {
    throw std::invalid_argument("ImageProcessor::gaussianBlur: input image is empty");
  }
  if (kernel_size % 2 == 0) {
    kernel_size += 1;
  }
  cv::Mat blurred;
  cv::GaussianBlur(input, blurred, cv::Size(kernel_size, kernel_size), 0);
  return blurred;
}
\end{lstlisting}

\noindent\textbf{\texttt{convolutionFilter}}
\begin{lstlisting}
cv::Mat ImageProcessor::convolutionFilter(const cv::Mat & input) const
{
  if (input.empty()) {
    throw std::invalid_argument("ImageProcessor::convolutionFilter: input image is empty");
  }
  if (this->kernel_.empty()) {
    throw std::runtime_error("kernel is empty");
  }
  cv::Mat filtered;

  // rotate kernel 180 degrees so filter2D's correlation becomes convolution
  cv::rotate(kernel_, kernel_, cv::ROTATE_180);

  cv::filter2D(input, filtered, -1, kernel_, cv::Point(-1, -1), 0, cv::BORDER_CONSTANT);

  return filtered;
}
\end{lstlisting}

\noindent\textbf{\texttt{convolutionFilterManual}}
\begin{lstlisting}
cv::Mat ImageProcessor::convolutionFilterManual(const cv::Mat & input) const
{
  if (input.empty()) {
    throw std::invalid_argument("ImageProcessor::convolutionFilterManual: input image is empty");
  }
  if (kernel_.empty()) {
    throw std::runtime_error("kernel is empty");
  }
  if (input.depth() != CV_8U) {
    throw std::invalid_argument(
      "ImageProcessor::convolutionFilterManual: only 8-bit images are supported");
  }

  cv::Mat kernel_f;
  kernel_.convertTo(kernel_f, CV_32F);

  const int kh = kernel_f.rows;
  const int kw = kernel_f.cols;
  const int anchor_y = kh / 2;
  const int anchor_x = kw / 2;

  const int rows = input.rows;
  const int cols = input.cols;
  const int channels = input.channels();

  cv::Mat output = cv::Mat::zeros(input.size(), input.type());

  for (int y = 0; y < rows; ++y) {
    uchar * out_row = output.ptr<uchar>(y);

    for (int x = 0; x < cols; ++x) {
      for (int c = 0; c < channels; ++c) {
        float sum = 0.0f;

        for (int i = 0; i < kh; ++i) {
          int sy = y + i - anchor_y;
          if (sy < 0 || sy >= rows) {
            continue;  // zero padding
          }
          const uchar * in_row = input.ptr<uchar>(sy);

          for (int j = 0; j < kw; ++j) {
            int sx = x + j - anchor_x;
            if (sx < 0 || sx >= cols) {
              continue;  // zero padding
            }
            float weight = kernel_f.at<float>(kh - 1 - i, kw - 1 - j);
            sum += weight * static_cast<float>(in_row[sx * channels + c]);
          }
        }

        out_row[x * channels + c] = cv::saturate_cast<uchar>(sum);
      }
    }
  }

  return output;
}
\end{lstlisting}

\noindent\textbf{\texttt{fourierTransform}}
\begin{lstlisting}
cv::Mat ImageProcessor::fourierTransform(const cv::Mat & input, bool shift) const
{
  if (input.empty()) {
    throw std::invalid_argument("ImageProcessor::fourierTransform: input image is empty");
  }

  cv::Mat gray = toGrayscale(input);

  int opt_rows = cv::getOptimalDFTSize(gray.rows);
  int opt_cols = cv::getOptimalDFTSize(gray.cols);
  cv::Mat padded;
  cv::copyMakeBorder(
    gray, padded, 0, opt_rows - gray.rows, 0, opt_cols - gray.cols,
    cv::BORDER_CONSTANT, cv::Scalar::all(0));

  cv::Mat planes[] = {cv::Mat_<float>(padded), cv::Mat::zeros(padded.size(), CV_32F)};
  cv::Mat complex_img;
  cv::merge(planes, 2, complex_img);

  cv::dft(complex_img, complex_img);

  cv::split(complex_img, planes);
  cv::magnitude(planes[0], planes[1], planes[0]);
  cv::Mat magnitude = planes[0];

  magnitude += cv::Scalar::all(1);
  cv::log(magnitude, magnitude);

  magnitude = magnitude(cv::Rect(0, 0, magnitude.cols & -2, magnitude.rows & -2));

  if (shift) {
    int cx = magnitude.cols / 2;
    int cy = magnitude.rows / 2;

    cv::Mat q0(magnitude, cv::Rect(0, 0, cx, cy));
    cv::Mat q1(magnitude, cv::Rect(cx, 0, cx, cy));
    cv::Mat q2(magnitude, cv::Rect(0, cy, cx, cy));
    cv::Mat q3(magnitude, cv::Rect(cx, cy, cx, cy));

    cv::Mat tmp;
    q0.copyTo(tmp);
    q3.copyTo(q0);
    tmp.copyTo(q3);

    q1.copyTo(tmp);
    q2.copyTo(q1);
    tmp.copyTo(q2);
  }

  cv::normalize(magnitude, magnitude, 0, 255, cv::NORM_MINMAX);
  magnitude.convertTo(magnitude, CV_8U);

  return magnitude;
}
\end{lstlisting}

\twocolumn
\begin{thebibliography}{0}
\bibitem{opencv-filter2d} OpenCV Development Team, ``Making your own linear filters!,'' \textit{OpenCV Tutorials}. [Online]. Available: \url{https://docs.opencv.org/4.x/d4/dbd/tutorial_filter_2d.html}
\bibitem{opencv-dft} OpenCV Development Team, ``Discrete Fourier Transform,'' \textit{OpenCV Tutorials}. [Online]. Available: \url{https://docs.opencv.org/4.x/d8/d01/tutorial_discrete_fourier_transform.html}
\end{thebibliography}

\end{document}
